\documentclass[aspectratio=169,10pt]{beamer}
% ═══════════════════════════════════════════════════════════════════════════════
% SHARED PREAMBLE — Foundation Models Course (HIT)
% To be \input by each lecture file AFTER \documentclass
% ═══════════════════════════════════════════════════════════════════════════════

% ─────────────────────────────────────────────────────────────────────────────
% BEAMER THEME & BASIC SETUP
% ─────────────────────────────────────────────────────────────────────────────
\usetheme{Madrid}
\usecolortheme{default}

% Remove navigation symbols
\beamertemplatenavigationsymbolsempty

% ─────────────────────────────────────────────────────────────────────────────
% COLORS — HIT Branding
% ─────────────────────────────────────────────────────────────────────────────
\definecolor{hitblue}{RGB}{0,51,102}        % Primary: HIT Navy
\definecolor{hitgold}{RGB}{192,148,0}       % Accent: HIT Gold
\definecolor{hitlight}{RGB}{230,240,250}    % Light background

% Override Beamer palette
\setbeamercolor{primary}{fg=white,bg=hitblue}
\setbeamercolor{secondary}{fg=hitblue,bg=hitlight}
\setbeamercolor{structure}{fg=hitblue}
\setbeamercolor{alerted text}{fg=hitgold}

% ─────────────────────────────────────────────────────────────────────────────
% PACKAGES
% ─────────────────────────────────────────────────────────────────────────────
\usepackage{amsmath}
\usepackage{amssymb}
\usepackage{mathtools}
\usepackage{booktabs}
\usepackage{graphicx}
\usepackage{hyperref}
\usepackage{tikz}
\usepackage{pgfplots}
\usepackage{tcolorbox}
\usepackage{listings}
\usepackage{xcolor}
\usepackage{algorithm}
\usepackage{algpseudocode}

% ─────────────────────────────────────────────────────────────────────────────
% TikZ LIBRARIES
% ─────────────────────────────────────────────────────────────────────────────
\usetikzlibrary{shapes,arrows,positioning,calc,chains,scopes}
\pgfplotsset{compat=1.17}

% ─────────────────────────────────────────────────────────────────────────────
% CODE LISTINGS
% ─────────────────────────────────────────────────────────────────────────────
\lstset{
  basicstyle=\ttfamily\small,
  breaklines=true,
  commentstyle=\color{gray},
  keywordstyle=\color{hitblue},
  stringstyle=\color{hitgold},
  showstringspaces=false,
  backgroundcolor=\color{hitlight},
  frame=single,
  rulecolor=\color{hitblue}
}

% ─────────────────────────────────────────────────────────────────────────────
% TCOLORBOX STYLES
% ─────────────────────────────────────────────────────────────────────────────
\tcbset{
  hitbox/.style={
    colback=hitlight,
    colframe=hitblue,
    fonttitle=\bfseries,
    boxrule=1pt
  },
  alertbox/.style={
    colback=yellow!10,
    colframe=hitgold,
    fonttitle=\bfseries,
    boxrule=1.5pt
  }
}

% ─────────────────────────────────────────────────────────────────────────────
% CUSTOM COMMANDS
% ─────────────────────────────────────────────────────────────────────────────

% Placeholder for figures
\newcommand{\placeholder}[3]{
  \begin{center}
    \fbox{\begin{minipage}[c][#2][c]{#1}
      \centering\small\color{gray}[Figure: #3]
    \end{minipage}}
  \end{center}
}

% Section divider frame
\newcommand{\sectionframe}[2]{
  \begin{frame}[plain]
    \begin{tikzpicture}[remember picture,overlay]
      \fill[hitblue] (current page.south west) rectangle (current page.north east);
      \node[anchor=center, white, font=\Large\bfseries] at (current page.center) {
        \begin{tabular}{c}
          #1 \\[0.3cm]
          {\small\color{hitgold}#2}
        \end{tabular}
      };
    \end{tikzpicture}
  \end{frame}
}

% ─────────────────────────────────────────────────────────────────────────────
% LOAD CUSTOM MACROS
% ─────────────────────────────────────────────────────────────────────────────
% ═══════════════════════════════════════════════════════════════════════════════
% SHARED MACROS — Mathematical Notation for Foundation Models Course
% ═══════════════════════════════════════════════════════════════════════════════

% ─────────────────────────────────────────────────────────────────────────────
% ACTIVATION & LAYER FUNCTIONS
% ─────────────────────────────────────────────────────────────────────────────
\newcommand{\softmax}{\mathrm{softmax}}
\newcommand{\sigmoid}{\mathrm{sigmoid}}
\newcommand{\relu}{\mathrm{ReLU}}
\newcommand{\gelu}{\mathrm{GELU}}
\newcommand{\LayerNorm}{\mathrm{LayerNorm}}
\newcommand{\FFN}{\mathrm{FFN}}
\newcommand{\MHA}{\mathrm{MHA}}
\newcommand{\Attn}{\mathrm{Attention}}

% ─────────────────────────────────────────────────────────────────────────────
% ATTENTION COMPONENTS
% ─────────────────────────────────────────────────────────────────────────────
\newcommand{\query}{Q}
\newcommand{\key}{K}
\newcommand{\val}{V}
\newcommand{\dmodel}{d_{\text{model}}}
\newcommand{\dk}{d_k}
\newcommand{\nheads}{h}

% Scaled dot-product attention formula
\newcommand{\SDPA}{
  \Attn(\query, \key, \val) = \softmax\left(\frac{\query \key^T}{\sqrt{\dk}}\right) \val
}

\newcommand{\CrossAttn}{\text{CrossAttention}}

% ─────────────────────────────────────────────────────────────────────────────
% PROBABILITY & EXPECTATION
% ─────────────────────────────────────────────────────────────────────────────
\newcommand{\E}{\mathbb{E}}
\newcommand{\Prob}{\mathbb{P}}
\newcommand{\KL}{\mathrm{KL}}
\newcommand{\ELBO}{\mathrm{ELBO}}
\newcommand{\given}{\mid}

% ─────────────────────────────────────────────────────────────────────────────
% NORMS & OPERATIONS
% ─────────────────────────────────────────────────────────────────────────────
\newcommand{\norm}[1]{\left\lVert #1 \right\rVert}
\newcommand{\abs}[1]{\left| #1 \right|}
\newcommand{\inner}[2]{\langle #1, #2 \rangle}
\newcommand{\T}{^{\top}}

% ─────────────────────────────────────────────────────────────────────────────
% VECTORS & MATRICES
% ─────────────────────────────────────────────────────────────────────────────
\newcommand{\bvec}[1]{\mathbf{#1}}
\newcommand{\mat}[1]{\mathbf{#1}}
\newcommand{\iid}{\stackrel{\text{i.i.d.}}{\sim}}
\newcommand{\defeq}{\coloneq}

% ─────────────────────────────────────────────────────────────────────────────
% LANGUAGE MODELING
% ─────────────────────────────────────────────────────────────────────────────
\newcommand{\vocab}{\mathcal{V}}
\newcommand{\context}{\text{context}}
\newcommand{\token}{t}
\newcommand{\seq}{x}
\newcommand{\ptheta}{p_\theta}
\newcommand{\pref}{p_{\text{ref}}}

% ─────────────────────────────────────────────────────────────────────────────
% RLHF & DPO
% ─────────────────────────────────────────────────────────────────────────────
\newcommand{\piref}{\pi_{\text{ref}}}
\newcommand{\pitheta}{\pi_\theta}
\newcommand{\yw}{y_w}
\newcommand{\yl}{y_l}
\newcommand{\DPOloss}{\mathcal{L}_{\mathrm{DPO}}}

% ─────────────────────────────────────────────────────────────────────────────
% RETRIEVAL & RAG
% ─────────────────────────────────────────────────────────────────────────────
\newcommand{\docs}{\mathcal{D}}
\newcommand{\qvec}{q}
\newcommand{\dvec}{d}

% ─────────────────────────────────────────────────────────────────────────────
% AGENTS & REASONING
% ─────────────────────────────────────────────────────────────────────────────
\newcommand{\obs}{o}
\newcommand{\act}{a}
\newcommand{\thought}{t}
\newcommand{\tools}{\mathcal{T}}
\newcommand{\history}{h}
\newcommand{\concat}{\oplus}

% ─────────────────────────────────────────────────────────────────────────────
% SETS & LOGIC
% ─────────────────────────────────────────────────────────────────────────────
\newcommand{\R}{\mathbb{R}}
\newcommand{\N}{\mathbb{N}}
\newcommand{\Z}{\mathbb{Z}}

\endinput


% ─────────────────────────────────────────────────────────────────────────────
% TITLE & AUTHOR (can be overridden in individual lectures)
% ─────────────────────────────────────────────────────────────────────────────
\author[D. Lederman]{Dror Lederman}
\institute{Holon Institute of Technology (HIT) \\ Data Science Department}
\date{\today}

% ─────────────────────────────────────────────────────────────────────────────
% FOOTER & HEADER
% ─────────────────────────────────────────────────────────────────────────────
\setbeamertemplate{footline}{
  \leavevmode%
  \hbox{%
    \begin{beamercolorbox}[wd=0.33\paperwidth,ht=2.25ex,dp=1ex,center]{author in head/foot}%
      \usebeamerfont{author in head/foot}\insertshortauthor
    \end{beamercolorbox}%
    \begin{beamercolorbox}[wd=0.34\paperwidth,ht=2.25ex,dp=1ex,center]{title in head/foot}%
      \usebeamerfont{title in head/foot}\insertshorttitle
    \end{beamercolorbox}%
    \begin{beamercolorbox}[wd=0.33\paperwidth,ht=2.25ex,dp=1ex,right]{frame number}%
      \usebeamerfont{frame number}\insertframenumber{} / \inserttotalframenumber\hspace*{2ex}
    \end{beamercolorbox}%
  }%
  \vskip0pt%
}

\endinput


\title{Week 5: Embeddings and Semantic Search}
\subtitle{Foundation Models and Agentic AI}
\author{Lecturer Name}
\institute{Holon Institute of Technology (HIT)}
\date{Semester, 2025}

\begin{document}

\begin{frame}
  \titlepage
\end{frame}

\begin{frame}{Learning Objectives}
  After this lecture you will be able to:
  \begin{itemize}
    \item Explain \textbf{representation learning} and what makes a good embedding.
    \item Derive \textbf{Sentence-BERT} and contrastive fine-tuning for sentence embeddings.
    \item Compare \textbf{dense retrieval} (DPR) and \textbf{sparse retrieval} (BM25).
    \item Understand \textbf{approximate nearest neighbor (ANN)} search algorithms (HNSW, IVF).
    \item Describe the architecture of \textbf{vector databases} and \textbf{hybrid search}.
    \item Explain \textbf{late interaction} (ColBERT) and when it outperforms bi-encoders.
  \end{itemize}
\end{frame}

\section{Representation Learning}
\sectionframe{Representation Learning}{What is a good embedding?}

\begin{frame}{What Is an Embedding?}
  \begin{tcolorbox}[hitbox, title=Definition]
    An \textbf{embedding} is a function $f: \mathcal{X} \to \R^d$ that maps inputs
    (tokens, sentences, images, documents) to a continuous vector space such that
    \textbf{semantic similarity correlates with geometric proximity}.
  \end{tcolorbox}
  \vspace{0.4em}
  \textbf{Desired properties:}
  \begin{itemize}
    \item Synonyms should be close: $f(\text{``car''}) \approx f(\text{``automobile''})$
    \item Distinct concepts should be far: $f(\text{``king''}) \not\approx f(\text{``peasant''})$
    \item Analogy structure: $f(\text{king}) - f(\text{man}) + f(\text{woman}) \approx f(\text{queen})$
  \end{itemize}
  \vspace{0.3em}
  \textbf{Applications:} Semantic search, clustering, classification, RAG retrieval.
\end{frame}

\begin{frame}{From Word2Vec to Contextualized Embeddings}
  \begin{center}
  \begin{tabular}{lll}
    \toprule
    \textbf{Model} & \textbf{Type} & \textbf{Key limitation} \\
    \midrule
    Word2Vec / GloVe & Static word & One vector per word (no context) \\
    ELMo & Contextual (LSTM) & Shallow fusion \\
    BERT [CLS] & Contextual (Transformer) & Not optimized for similarity \\
    Sentence-BERT & Sentence & Fine-tuned for semantic similarity \\
    E5, BGE, GTE & Modern dense & Fine-tuned on massive IR datasets \\
    \bottomrule
  \end{tabular}
  \end{center}
  \vspace{0.4em}
  \textbf{Key point:} Raw BERT [CLS] embeddings perform poorly for semantic search.
  Task-specific fine-tuning (or contrastive training) is necessary.
  \note{Humeau et al. (2019) showed that raw BERT is worse than GloVe averaging for semantic textual similarity.}
\end{frame}

\section{Sentence Embeddings}
\sectionframe{Sentence Embeddings}{Sentence-BERT and beyond}

\begin{frame}{Sentence-BERT (SBERT)}
  \textbf{Reimers and Gurevych (2019)}~\cite{reimers2019sentence}
  \vspace{0.4em}
  \begin{columns}[T]
    \column{0.55\textwidth}
      \textbf{Architecture:} Siamese BERT network.
      \begin{itemize}
        \item Two BERT encoders sharing weights.
        \item Mean-pool the token embeddings to get sentence vectors.
        \item Fine-tune with \textbf{Natural Language Inference (NLI)} data using contrastive loss.
      \end{itemize}
      \vspace{0.3em}
      \textbf{Loss (NLI triplet):}
      \[
        \mathcal{L} = \max(0, \norm{f(a) - f(p)} - \norm{f(a) - f(n)} + \epsilon)
      \]
    \column{0.41\textwidth}
      \placeholder{0.95\textwidth}{4cm}{SBERT: siamese encoding of sentence pairs}
  \end{columns}
  \note{The key win: SBERT reduces pairwise similarity computation from $O(n^2)$ cross-encoder forward passes to $O(n)$ encodings + a single inner product.}
\end{frame}

\begin{frame}{Modern Embedding Models}
  \begin{columns}[T]
    \column{0.52\textwidth}
      \textbf{Training recipe for state-of-the-art embeddings:}
      \begin{enumerate}
        \item Pretrain on contrastive web-scale data (MS-MARCO, CC-News pairs).
        \item Hard negative mining: retrieve top-$k$ near-miss negatives.
        \item Knowledge distillation from a cross-encoder teacher.
        \item Optional: instruction-tuned embeddings (E5, Gecko).
      \end{enumerate}
    \column{0.44\textwidth}
      \textbf{Current leaders (MTEB benchmark):}
      \begin{itemize}
        \item E5-mistral-7b
        \item BGE-M3 (multi-lingual)
        \item GTE-Qwen
        \item Cohere Embed v3
        \item OpenAI text-embedding-3-large
      \end{itemize}
  \end{columns}
\end{frame}

\section{Sparse vs.\ Dense Retrieval}
\sectionframe{Retrieval Paradigms}{BM25 vs.\ Dense Passage Retrieval}

\begin{frame}{BM25: Sparse Term-Based Retrieval}
  \begin{tcolorbox}[hitbox, title=BM25 Score]
    \[
      \text{BM25}(q, d) = \sum_{t \in q} \mathrm{IDF}(t)
      \cdot \frac{f(t,d)(k_1+1)}{f(t,d) + k_1\left(1 - b + b\cdot\frac{|d|}{\text{avgdl}}\right)}
    \]
  \end{tcolorbox}
  \vspace{0.3em}
  \begin{itemize}
    \item $f(t,d)$: term frequency in document $d$; $\mathrm{IDF}(t)$: inverse document frequency.
    \item $k_1 \approx 1.2$, $b \approx 0.75$: saturation and document-length normalization.
    \item Extremely fast: inverted index lookups.
    \item \textbf{Weakness:} vocabulary mismatch — misses synonyms, paraphrases.
  \end{itemize}
\end{frame}

\begin{frame}{Dense Passage Retrieval (DPR)}
  \textbf{Karpukhin et al.\ (2020)}~\cite{karpukhin2020dpr}
  \vspace{0.4em}
  \begin{columns}[T]
    \column{0.55\textwidth}
      \textbf{Bi-encoder setup:}
      \begin{itemize}
        \item Query encoder $E_Q$, passage encoder $E_P$ (both BERT).
        \item Retrieve top-$k$ by inner product: $E_Q(q)\T E_P(p)$.
        \item Fine-tuned on labeled QA pairs with in-batch negatives.
      \end{itemize}
      \vspace{0.4em}
      \textbf{Advantages:}
      \begin{itemize}
        \item Captures semantic similarity, not just term overlap.
        \item Passage embeddings precomputed and cached.
      \end{itemize}
    \column{0.41\textwidth}
      \placeholder{0.95\textwidth}{3.5cm}{DPR: query encoder vs.\ passage encoder, ANN retrieval}
  \end{columns}
\end{frame}

\begin{frame}{Hybrid Search}
  \begin{tcolorbox}[hitbox, title=Hybrid Search = BM25 + Dense]
    Combine sparse and dense scores (e.g., Reciprocal Rank Fusion):
    \[
      \text{RRF}(d) = \sum_{r \in \{r_\text{sparse}, r_\text{dense}\}} \frac{1}{k + r(d)}
    \]
  \end{tcolorbox}
  \vspace{0.4em}
  \begin{itemize}
    \item Sparse: exact keyword match (good for rare entities, product codes, UUIDs).
    \item Dense: semantic similarity (good for paraphrases, concept queries).
    \item Hybrid consistently outperforms either alone on BEIR benchmark.
    \item Used by: Elasticsearch, Weaviate, Pinecone, Azure Cognitive Search.
  \end{itemize}
\end{frame}

\section{Approximate Nearest Neighbor Search}
\sectionframe{ANN Search}{Scaling to billions of vectors}

\begin{frame}{The ANN Problem}
  \begin{tcolorbox}[alertbox, title=Problem]
    Exact nearest neighbor search in $\R^d$ requires $O(n \cdot d)$ per query — intractable for
    $n = 10^9$ vectors at $d = 768$.
  \end{tcolorbox}
  \vspace{0.4em}
  \textbf{ANN approaches:}
  \begin{enumerate}
    \item \textbf{LSH} (Locality-Sensitive Hashing): hash similar vectors to same bucket.
    \item \textbf{IVF} (Inverted File Index): cluster vectors, search only nearby clusters.
    \item \textbf{HNSW} (Hierarchical Navigable Small Worlds)~\cite{malkov2018hnsw}: hierarchical graph, greedy search.
    \item \textbf{Product Quantization (PQ)}: compress vectors, reduce memory 8–32×.
  \end{enumerate}
  \vspace{0.3em}
  HNSW is the most widely adopted: $O(\log n)$ query time, high recall.
\end{frame}

\begin{frame}{Vector Databases}
  \begin{columns}[T]
    \column{0.50\textwidth}
      \textbf{Popular options:}
      \begin{itemize}
        \item \textbf{FAISS} (Meta): library, not a DB; supports IVF + PQ + HNSW.
        \item \textbf{Pinecone}: managed cloud, serverless, hybrid search.
        \item \textbf{Weaviate}: open-source, schema-based, multi-modal.
        \item \textbf{Qdrant}: open-source, payload filtering, Rust.
        \item \textbf{Chroma}: lightweight, in-memory, for prototyping.
        \item \textbf{pgvector}: Postgres extension; no separate infra.
      \end{itemize}
    \column{0.46\textwidth}
      \textbf{Key features to compare:}
      \begin{itemize}
        \item Indexing algorithm (HNSW vs.\ IVF)
        \item Metadata filtering (pre-filter vs.\ post-filter)
        \item Persistent vs.\ in-memory
        \item Multi-vector support (ColBERT-style)
        \item Replication and sharding
      \end{itemize}
  \end{columns}
\end{frame}

\section{Late Interaction: ColBERT}
\sectionframe{ColBERT}{Beyond single-vector representations}

\begin{frame}{ColBERT: Contextualized Late Interaction}
  \textbf{Khattab and Zaharia (2020)}~\cite{khattab2020colbert}
  \vspace{0.4em}
  \begin{tcolorbox}[hitbox, title=Late Interaction Scoring]
    Represent query $q$ and document $d$ as \emph{sequences} of token embeddings:
    \[
      \text{Score}(q, d) = \sum_{i \in q} \max_{j \in d} E_q^i \cdot E_d^j
    \]
    MaxSim over document tokens for each query token; sum over query.
  \end{tcolorbox}
  \vspace{0.3em}
  \begin{itemize}
    \item Better than bi-encoder (less information compression).
    \item Much faster than cross-encoder (offline document encoding).
    \item Used in: ColBERTv2, PLAID retrieval pipeline, RAGatouille.
  \end{itemize}
\end{frame}

\section{Paper Discussion}
\begin{frame}{Paper Discussion}
  \papercite{Dense Passage Retrieval for Open-Domain Question Answering}
    {Karpukhin, Oguz, Min, Lewis et al.\ (Facebook AI)}
    {EMNLP 2020~\cite{karpukhin2020dpr}}
  \vspace{0.4em}
  \textbf{Key findings:}
  \begin{itemize}
    \item DPR beats BM25 by $+9$ points on Natural Questions top-20 retrieval accuracy.
    \item In-batch negatives are a simple but effective training strategy.
    \item Showed that \emph{dense retrieval from scratch} can match BM25 with minimal data.
  \end{itemize}
  \vspace{0.3em}
  \textbf{Discussion:} DPR requires pre-encoding all passages. What are the engineering challenges
  of maintaining a dense index as the document corpus updates in real time?
\end{frame}

\section{Lab / Demo}
\begin{frame}{Lab Idea: Build a Semantic Search Engine}
  \begin{tcolorbox}[hitbox, title=Lab 5 — Semantic Search with FAISS]
    \textbf{Goal:} Build and evaluate a working semantic search system.\\[4pt]
    \textbf{Tools:} \texttt{sentence-transformers}, \texttt{faiss-cpu}, Wikipedia dump (10K articles)
    \begin{enumerate}
      \item Encode 10K Wikipedia passages with \texttt{all-MiniLM-L6-v2}.
      \item Build a FAISS flat index; measure exact search latency.
      \item Switch to FAISS HNSW; compare speed/recall trade-off.
      \item Implement BM25 (pyserini); compare result quality with 10 queries.
      \item Implement RRF hybrid search; measure MRR@10 on TriviaQA.
    \end{enumerate}
    \textbf{Deliverable:} Speed/recall table + qualitative comparison.
  \end{tcolorbox}
\end{frame}

\section{Discussion \& Summary}
\begin{frame}{Discussion Questions}
  \begin{enumerate}
    \item BM25 has been around since the 1990s. Why does it still often outperform dense retrieval on out-of-domain benchmarks?
    \item HNSW is the dominant ANN algorithm, but it has no formal retrieval guarantees. What are the real-world failure modes?
    \item ColBERT stores per-token embeddings for every document. How does this change the storage/compute trade-off compared to a bi-encoder?
    \item When would you use a managed vector database (Pinecone) vs.\ running your own (Qdrant/FAISS)? What are the key decision factors?
  \end{enumerate}
\end{frame}

\begin{frame}{Summary}
  \begin{itemize}
    \item Good embeddings require \textbf{contrastive fine-tuning}; raw BERT performs poorly for search.
    \item \textbf{Sentence-BERT} and modern models (E5, BGE) encode sentences into semantically meaningful spaces.
    \item \textbf{Sparse retrieval} (BM25): fast, exact, vocabulary-limited.
    \item \textbf{Dense retrieval} (DPR): slower offline encoding, better semantic matching.
    \item \textbf{Hybrid search} (BM25 + dense) consistently outperforms either alone.
    \item \textbf{HNSW} enables sub-linear ANN search; \textbf{ColBERT} improves quality via late interaction.
  \end{itemize}
  \vspace{0.3em}
  \textbf{Next week:} Retrieval-Augmented Generation — putting search inside a generative model.
\end{frame}

\begin{frame}[allowframebreaks]{Suggested Reading}
  \nocite{reimers2019sentence, karpukhin2020dpr, khattab2020colbert, malkov2018hnsw}
  \bibliography{../../references}
\end{frame}

\end{document}
